\chapter{Aus Welt und Zeit}
\quellseite{9}

Als ein Berauschungsmittel wird mehr und mehr, namentlich in den Südstaaten unsers Landes, Cocain gebraucht. Ein weltliches Blatt berichtet: In Atlanta, Ga., hat das Einathmen von Cocain einen solchen Umfang angenommen, daß Schankwirthe, deren Kunden größtentheils Farbige waren, gezwungen sind, ihre Locale zu schließen; denn wenn die Gewohnheit des Cocaingenusses in einem Menschen Wurzel faßt, begehrt er nicht länger alkoholische Getränke. Cocain ist billiger als Branntwein, und die Wirkung dauert länger. Für 25 oder 50 Cents kann man einen Wochenvorrath Cocain kaufen. Die Cocainkrystalle werden pulverisirt und mit Zucker vermischt, dann in eine Schachtel gelegt und in den Taschen getragen, so daß jeden Augenblick eine Dosis genommen werden kann. Durch Einathmen des Cocains durch die Nasenlöcher bringt es in das Gehirn rascher, als wenn es verschluckt oder eingespritzt wird. Die Wirkungen des Cocains haben Aehnlichkeit mit denen des Morphiums. Wie Opium erzeugt es Visionen von Reichthum, Schönheit und Glückseligkeit. Unzweifelhaft greift Cocain das Gehirn rasch an, und die Folge ist, daß die Irrenanstalten des Südens von den unglücklichen Opfern dieses berauschenden Giftes bezogen werden. — Die Sünde ist der Leute Verderben.

Ueber die sittliche Fäulniß in Paris, schreibt die „Allgemeine Ev.-Luth. Kirchenzeitung“, liegen Nachrichten vor, die geradezu mit Entsetzen erfüllen. Die Schamlosigkeit hat auch die Straße in einer Weise in Besitz genommen, von der man anderswo schwerlich eine Vorstellung hat. Kaum einen Schritt kann man auf der Straße gehen, ohne ausschweifende Bilder, farbige Anschläge, illustrirte Schriften, Zeichnungen, Photographien, Transparenkarten in unzähligen Verschiedenheiten des Wollust und des Rothes zu erblicken. Das Scheußlichste ist, daß nicht einmal die Gewinnsucht einen Erklärungsgrund abgibt; vielmehr werden Hunderttausende von infamen Veröffentlichungen ohne Gewinn und Nutzen ins Publicum geworfen und besonders verschwenderisch in die Jugend. So wurden vor Kurzem an Schüler und junge Mädchen auf den Haltestellen der Straßenbahn kleine Broschüren vertheilt mit der Ueberschrift: „Für Gott“ oder „Für’s Vaterland“, die man vertrauensvoll annahm und die nichts enthielten als unbeschreibbaren Schmutz. Irgend ein materieller Gewinn ist also hier gar nicht ins Auge gefaßt; es ist die Lust an der Verführung, der Satanismus in vollendeter Gestalt“. Ein Seitenstück dazu sind die schamlosen Abbildungen, die die Leiter gemeiner Theater in allen Städten unsers Landes öffentlich aufhängen dürfen, und die schon den Kindern, die fast nicht umhin können, sie zu sehen, sittliche Gefahren bringen.
